% 本书原创：同一欧氏比例绘出非正交格与逆转置对偶格。
\begin{figure}[htbp]
\centering
\begin{tikzpicture}[x=1cm,y=1cm,>=stealth,every node/.style={font=\small},line width=.65pt]
\begin{scope}[xshift=0.0cm]
\fill[AnalysisBlue!10] (0.00000,0.00000)--(1.50000,0.00000)--(2.00000,1.00000)--(0.50000,1.00000)--cycle;
\draw[->,gray] (-1.2,0)--(3.4,0);\draw[->,gray] (0,-1.2)--(0,2.6);
\fill[gray] (-1.00000,1.00000) circle(1.15pt);
\fill[gray] (-0.50000,2.00000) circle(1.15pt);
\fill[gray] (-0.50000,-1.00000) circle(1.15pt);
\fill[gray] (0.00000,0.00000) circle(1.15pt);
\fill[gray] (0.50000,1.00000) circle(1.15pt);
\fill[gray] (1.00000,2.00000) circle(1.15pt);
\fill[gray] (1.00000,-1.00000) circle(1.15pt);
\fill[gray] (1.50000,0.00000) circle(1.15pt);
\fill[gray] (2.00000,1.00000) circle(1.15pt);
\fill[gray] (2.50000,2.00000) circle(1.15pt);
\fill[gray] (2.50000,-1.00000) circle(1.15pt);
\fill[gray] (3.00000,0.00000) circle(1.15pt);
\draw[->,AnalysisBlue,thick] (0,0)--(1.50000,0.00000) node[above right] {\(b_1\)};
\draw[->,AnalysisBlue,thick] (0,0)--(0.50000,1.00000) node[above right] {\(b_2\)};
\node at (1,3.0) {\(\Lambda,\quad v(\Lambda)=3/2\)};\end{scope}
\begin{scope}[xshift=6.6cm]
\fill[AnalysisBlue!10] (0.00000,0.00000)--(0.66667,-0.33333)--(0.66667,0.66667)--(0.00000,1.00000)--cycle;
\draw[->,gray] (-1.2,0)--(3.4,0);\draw[->,gray] (0,-1.2)--(0,2.6);
\fill[gray] (-0.66667,-0.66667) circle(1.15pt);
\fill[gray] (-0.66667,0.33333) circle(1.15pt);
\fill[gray] (-0.66667,1.33333) circle(1.15pt);
\fill[gray] (0.00000,-1.00000) circle(1.15pt);
\fill[gray] (0.00000,0.00000) circle(1.15pt);
\fill[gray] (0.00000,1.00000) circle(1.15pt);
\fill[gray] (0.00000,2.00000) circle(1.15pt);
\fill[gray] (0.66667,-0.33333) circle(1.15pt);
\fill[gray] (0.66667,0.66667) circle(1.15pt);
\fill[gray] (0.66667,1.66667) circle(1.15pt);
\fill[gray] (1.33333,-0.66667) circle(1.15pt);
\fill[gray] (1.33333,0.33333) circle(1.15pt);
\fill[gray] (1.33333,1.33333) circle(1.15pt);
\fill[gray] (2.00000,-1.00000) circle(1.15pt);
\fill[gray] (2.00000,0.00000) circle(1.15pt);
\fill[gray] (2.00000,1.00000) circle(1.15pt);
\fill[gray] (2.00000,2.00000) circle(1.15pt);
\fill[gray] (2.66667,-0.33333) circle(1.15pt);
\fill[gray] (2.66667,0.66667) circle(1.15pt);
\fill[gray] (2.66667,1.66667) circle(1.15pt);
\draw[->,AnalysisBlue,thick] (0,0)--(0.66667,-0.33333) node[above right] {\(b_1^*\)};
\draw[->,AnalysisBlue,thick] (0,0)--(0.00000,1.00000) node[above right] {\(b_2^*\)};
\node at (1,3.0) {\(\Lambda^*,\quad v(\Lambda^*)=2/3\)};\end{scope}
\end{tikzpicture}
\caption[非正交格与对偶格]{非正交格与对偶格。左侧基矩阵为 \(\left(\begin{smallmatrix}3/2&1/2\\0&1\end{smallmatrix}\right)\)，右侧取逆转置；淡蓝色区域为相应基本平行四边形。两图采用同一长度单位，区域面积互为倒数。}
\label{b2:fig:06-dual-lattice}
\end{figure}
